\uuid{JKRX}
\exo7id{1575}
\titre{exo7 1575}
\auteur{legall}
\organisation{exo7}
\datecreate{1998-09-01}
\isIndication{false}
\isCorrection{false}
\chapitre{Réduction d'endomorphisme, polynôme annulateur}
\sousChapitre{Polynôme annulateur}
\cpge{mpsi-mp2i.espaces-vectoriels.a, pcsi.espaces-vectoriels.a, ptsi.geometrie-espace.c, mpsi-mp2i.applications-lineaires.d, pcsi.applications-lineaires.d, ptsi.applications-lineaires.d, mp-mpi.reduction-endomorphismes.a, pc.complements-espaces-vectoriels.a, psi.complements-espaces-vectoriels.a, pt.complements-algebre-lineaire.a, mp-mpi.reduction-endomorphismes.g, pc.complements-espaces-vectoriels.d, psi.complements-espaces-vectoriels.d}
\module{Algèbre}
\niveau{L2}
\difficulte{}

\contenu{
\question{
Soit $  P\in \C [X]  $ tel que $  P(0)=0  $ et $  P'(0)\not = 0  .$
Soit $  E  $ un $  \C $-espace vectoriel de dimension finie et $  f \in \mathcal{L} (E)  $ telle
que $  P(f)=0  .$ 

Montrer que $  \hbox{Ker}(f)=\hbox{Ker}(f^2)  ;$ en d\'eduire $  E=\hbox{Ker}(f)\oplus
\hbox{Im}(f).$
}
}
